\documentclass[aspectratio=169,11pt]{beamer}

\usetheme{Madrid}
\usecolortheme{seahorse}
\setbeamertemplate{navigation symbols}{}
\setbeamertemplate{footline}[frame number]

\usepackage[utf8]{inputenc}
\usepackage[T1]{fontenc}
\usepackage[spanish,es-tabla]{babel}
\usepackage{booktabs}
\usepackage{array}
\usepackage{amssymb}
\usepackage{amsmath}
\usepackage{tikz}
\usetikzlibrary{arrows.meta,positioning,shapes.geometric}

\title{ME\_3: Números Índices}
\subtitle{Cambio relativo, comparación temporal e índices económicos}
\author{Material de apoyo para estudio}
\date{\today}

\newcommand{\pregunta}[1]{\textbf{Pregunta que responde:} \emph{#1}}
\newcommand{\errorcomun}[1]{\textbf{Error común:} #1}
\newcommand{\diferencia}[1]{\textbf{Diferencia clave:} #1}
\newcommand{\intuicion}[1]{\textbf{Intuición:} #1}
\newcommand{\base}{\text{base}}
\newcommand{\actual}{\text{actual}}

\begin{document}

\begin{frame}
  \titlepage
\end{frame}

\begin{frame}{Cómo usar esta unidad}
\begin{itemize}
  \item Esta unidad continúa el bloque de estadística descriptiva: ahora el foco está en comparar magnitudes entre períodos.
  \item La estructura se mantiene: definición, intuición, pregunta, diferencia, error común y ejemplo.
  \item Los ejercicios típicos combinan precios, cantidades, valores, ponderaciones y elección de año base.
\end{itemize}
\end{frame}

\begin{frame}{Mapa de la Unidad 3}
\small
\begin{columns}[T]
\column{0.48\textwidth}
\begin{enumerate}
  \item Introducción a números índices
  \item Índices simples: precio, cantidad y valor
  \item Índices complejos no ponderados
  \item Índices económicos clásicos
  \item Índices ponderados
\end{enumerate}
\column{0.48\textwidth}
\begin{enumerate}\setcounter{enumi}{5}
  \item Laspeyres
  \item Paasche
  \item Relación precio--cantidad--valor
  \item Fisher
  \item Errores comunes
\end{enumerate}
\end{columns}
\end{frame}

\section{Introducción a los números índices}

\begin{frame}{¿Qué es un número índice?}
\textbf{Definición formal:} razón que compara el nivel de una variable en un período con su nivel en un período base.
\[
I_t=\frac{x_t}{x_0}
\]
\begin{itemize}
  \item \intuicion{mide cuánto representa el valor actual respecto del valor de referencia.}
  \item \pregunta{¿Cuánto cambió una magnitud en relación con el período base?}
  \item Es una razón adimensional: las unidades se cancelan.
  \item \textbf{Ejemplo:} si un precio sube de \$500 a \$650, entonces $I=650/500=1{,}30$.
  \item \errorcomun{confundir el índice $1{,}30$ con un aumento de $1{,}30\%$; el aumento es $30\%$.}
\end{itemize}
\end{frame}

\begin{frame}{Interpretación básica de un índice}
\begin{center}
\begin{tabular}{c p{0.58\textwidth} c}
\toprule
Índice & Interpretación & Variación \\
\midrule
$I=1$ & el nivel actual coincide con el período base & $0\%$ \\
$I>1$ & el nivel actual es mayor que el período base & positiva \\
$I<1$ & el nivel actual es menor que el período base & negativa \\
\bottomrule
\end{tabular}
\end{center}
\begin{block}{Conversión práctica}
Un índice puede expresarse como base 1 o base 100: $1{,}25=125$ si se usa base 100.
\end{block}
\diferencia{el índice compara niveles; la variación porcentual expresa el cambio neto respecto de la base.}
\end{frame}

\begin{frame}{Año base o período base}
\textbf{Definición formal:} período usado como referencia para comparar todos los demás períodos.
\begin{itemize}
  \item \intuicion{es el ``punto cero'' de la comparación temporal.}
  \item \pregunta{¿Respecto de qué momento estoy midiendo el cambio?}
  \item Conviene elegir un período normal, representativo y con datos confiables.
  \item Mantener una referencia fija permite comparar varios años entre sí mediante el mismo patrón.
  \item \errorcomun{cambiar el año base sin avisar y comparar índices que ya no tienen la misma referencia.}
\end{itemize}
\end{frame}

\begin{frame}{Variación porcentual a partir de un índice}
\[
\%\Delta=(I-1)\cdot 100
\]
\begin{itemize}
  \item Si $I=1{,}18$, entonces $\%\Delta=(1{,}18-1)100=18\%$.
  \item Si $I=0{,}92$, entonces $\%\Delta=(0{,}92-1)100=-8\%$.
  \item En base 100, la variación se obtiene como $I_{100}-100$.
  \item \pregunta{¿El nivel actual aumentó o disminuyó, y en qué porcentaje?}
  \item \errorcomun{decir que $I=92$ significa aumento de $92\%$; en base 100 significa caída de $8\%$.}
\end{itemize}
\end{frame}

\section{Números índices simples}

\begin{frame}{Concepto de índice simple}
\textbf{Definición formal:} índice que estudia la evolución de una sola variable o magnitud entre dos períodos.
\[
I(x)=\frac{x_t}{x_0}
\]
\begin{itemize}
  \item \intuicion{compara una cosa contra sí misma en otro momento.}
  \item \textbf{Magnitudes típicas:} precio, cantidad vendida, producción, ingreso o valor monetario.
  \item \diferencia{simple significa una variable o producto; complejo significa conjunto de productos o magnitudes.}
\end{itemize}
\end{frame}

\begin{frame}{Índice simple de precios}
\[
I(P)=\frac{P_t}{P_0}
\]
\begin{itemize}
  \item Mide la evolución temporal del precio de un bien.
  \item \pregunta{¿Cuánto cambió el precio respecto del período base?}
  \item \textbf{Ejemplo:} pan de \$1.000 a \$1.250: $I(P)=1{,}25$, aumento de $25\%$.
  \item \textbf{Interpretación económica:} aproxima inflación de ese producto específico.
  \item \errorcomun{generalizar la inflación de un producto a toda la economía.}
\end{itemize}
\end{frame}

\begin{frame}{Índice simple de cantidades}
\[
I(Q)=\frac{Q_t}{Q_0}
\]
\begin{itemize}
  \item Mide cambios en unidades físicas, producción, consumo o ventas.
  \item \pregunta{¿Se vendió o produjo más o menos que en la base?}
  \item \textbf{Ejemplo:} ventas de 80 a 100 unidades: $I(Q)=1{,}25$, aumento de $25\%$.
  \item \diferencia{un precio puede subir aunque la cantidad vendida baje; son dimensiones distintas.}
\end{itemize}
\end{frame}

\begin{frame}{Índice simple de valor}
\textbf{Valor:} monto monetario que combina precio y cantidad.
\[
V=P\cdot Q \qquad I(V)=\frac{V_t}{V_0}=\frac{P_tQ_t}{P_0Q_0}
\]
\begin{itemize}
  \item \intuicion{resume ingresos, ventas monetarias o gasto total.}
  \item Relación fundamental en índices simples:
  \[
  I(V)=I(P)\cdot I(Q)
  \]
  \item \textbf{Ejemplo:} si el precio sube $20\%$ y la cantidad baja $10\%$, entonces $I(V)=1{,}20\cdot0{,}90=1{,}08$.
  \item \errorcomun{interpretar un mayor valor como mayor cantidad vendida; puede deberse solo a precios.}
\end{itemize}
\end{frame}

\section{Números índices complejos}

\begin{frame}{Concepto de índice complejo}
\textbf{Definición formal:} índice que resume la evolución conjunta de varias magnitudes o productos.
\begin{itemize}
  \item \intuicion{pasa de mirar un producto a mirar una canasta.}
  \item \pregunta{¿Cómo cambió el conjunto completo de bienes o variables?}
  \item Se usa cuando interesa una visión global: canasta familiar, producción industrial, ventas de una empresa.
  \item \diferencia{el índice simple no requiere agregar productos; el complejo necesita un método de agregación.}
\end{itemize}
\end{frame}

\begin{frame}{Índices complejos no ponderados}
\textbf{Concepto:} todos los elementos reciben la misma importancia estadística.
\begin{itemize}
  \item \textbf{Ventaja:} cálculo rápido y fácil de explicar.
  \item \textbf{Limitación:} trata igual bienes muy distintos en relevancia económica.
  \item \textbf{Ejemplo:} dar el mismo peso a pan, arriendo y lápices puede distorsionar una canasta de consumo.
  \item \errorcomun{creer que no ponderar significa ser neutral; en realidad asigna igual peso a todo.}
\end{itemize}
\end{frame}

\begin{frame}{Media agregativa simple}
\[
I_A(P)=\frac{\sum P_t}{\sum P_0}
\]
\begin{itemize}
  \item Suma precios actuales y los compara con la suma de precios base.
  \item \pregunta{¿La suma total de magnitudes aumentó o disminuyó?}
  \item \textbf{Aplicación:} útil si las magnitudes son comparables y tienen escala similar.
  \item \errorcomun{sumar precios de bienes con unidades muy distintas sin justificar la agregación.}
\end{itemize}
\end{frame}

\begin{frame}{Método de los promedios}
\[
I_M(P)=\frac{1}{n}\sum_{i=1}^{n}\frac{P_{it}}{P_{i0}}
\]
\begin{itemize}
  \item Calcula primero índices individuales y luego los promedia.
  \item \pregunta{¿Cuál es el cambio relativo promedio de los productos?}
  \item \diferencia{la agregativa compara sumas; el método de promedios compara razones individuales.}
  \item Puede diferir mucho de la media agregativa cuando los precios base tienen escalas distintas.
\end{itemize}
\end{frame}

\begin{frame}{Comparación: agregativa simple vs promedio simple}
\small
\begin{center}
\begin{tabular}{lccccc}
\toprule
Producto & $P_0$ & $P_t$ & $P_t/P_0$ & Aporte a sumas & Aporte al promedio \\
\midrule
A & 100 & 120 & 1,20 & alto & igual \\
B & 10 & 20 & 2,00 & bajo & igual \\
\bottomrule
\end{tabular}
\end{center}
\[
I_A=\frac{120+20}{100+10}=1{,}27 \qquad
I_M=\frac{1{,}20+2{,}00}{2}=1{,}60
\]
\begin{block}{Idea clave}
Los métodos divergen cuando un producto pequeño cambia mucho o cuando las escalas iniciales son muy diferentes.
\end{block}
\end{frame}

\section{Índices económicos clásicos}

\begin{frame}{Índices económicos clásicos}
\begin{columns}[T]
\column{0.32\textwidth}
\textbf{Precios}
\begin{itemize}
  \item Inflación
  \item Poder adquisitivo
  \item Comparación temporal
\end{itemize}
\column{0.32\textwidth}
\textbf{Cantidades}
\begin{itemize}
  \item Producción
  \item Consumo
  \item Ventas físicas
\end{itemize}
\column{0.32\textwidth}
\textbf{Valor}
\begin{itemize}
  \item Ingresos
  \item Gasto total
  \item Efecto precio y cantidad
\end{itemize}
\end{columns}
\vspace{0.5em}
\errorcomun{concluir causalidad económica solo porque un índice aumentó o disminuyó.}
\end{frame}

\section{Índices compuestos ponderados}

\begin{frame}{Motivación de las ponderaciones}
\textbf{Idea central:} no todos los productos tienen la misma importancia en una canasta o en una economía.
\begin{itemize}
  \item Las ponderaciones representan participación, relevancia o peso económico.
  \item Evitan que un bien irrelevante tenga el mismo impacto que uno fundamental.
  \item \textbf{Ejemplo:} en el presupuesto familiar, el arriendo pesa más que un paquete de lápices.
  \item \pregunta{¿Qué productos deberían influir más en el índice global?}
\end{itemize}
\end{frame}

\begin{frame}{Concepto de ponderación}
\begin{itemize}
  \item \textbf{Peso estadístico:} coeficiente usado para dar importancia relativa a cada elemento.
  \item \textbf{Peso económico:} importancia medida por gasto, ingreso, producción o participación.
  \item En índices de precios, las cantidades suelen actuar como ponderaciones.
  \item En índices de cantidades, los precios suelen actuar como ponderaciones.
\end{itemize}
\begin{block}{Regla de lectura}
Precio se pondera con cantidad; cantidad se pondera con precio.
\end{block}
\end{frame}

\section{Índices de Laspeyres}

\begin{frame}{Índice de precios de Laspeyres}
\[
P_L=\frac{\sum P_tQ_0}{\sum P_0Q_0}
\]
\begin{itemize}
  \item Usa cantidades del período base como ponderaciones.
  \item \pregunta{¿Cuánto costaría hoy comprar la canasta base?}
  \item \textbf{Ventaja:} mantiene fija la canasta y facilita comparaciones.
  \item \textbf{Sesgo típico:} puede sobreestimar aumentos si los consumidores sustituyen bienes encarecidos. 
  \item \errorcomun{usar cantidades actuales en Laspeyres de precios.}
\end{itemize}
\end{frame}

\begin{frame}{Índice de cantidades de Laspeyres}
\[
Q_L=\frac{\sum Q_tP_0}{\sum Q_0P_0}
\]
\begin{itemize}
  \item Usa precios del período base como ponderaciones.
  \item \pregunta{¿Cómo cambió el volumen producido o vendido valorado a precios base?}
  \item \textbf{Aplicación:} comparar producción física evitando que los precios actuales distorsionen el cambio.
\end{itemize}
\end{frame}

\section{Índices de Paasche}

\begin{frame}{Índice de precios de Paasche}
\[
P_P=\frac{\sum P_tQ_t}{\sum P_0Q_t}
\]
\begin{itemize}
  \item Usa cantidades actuales como ponderaciones.
  \item \pregunta{¿Cuánto cuesta la canasta actual comparada a precios base?}
  \item Refleja patrones actuales de consumo o producción.
  \item \errorcomun{confundirlo con Laspeyres porque ambos comparan precios; la diferencia está en las cantidades usadas.}
\end{itemize}
\end{frame}

\begin{frame}{Índice de cantidades de Paasche}
\[
Q_P=\frac{\sum Q_tP_t}{\sum Q_0P_t}
\]
\begin{itemize}
  \item Usa precios actuales como ponderaciones.
  \item \pregunta{¿Cómo cambió el volumen usando la estructura de precios actual?}
  \item Puede incorporar cambios de importancia económica ocurridos en el período actual.
\end{itemize}
\end{frame}

\begin{frame}{Laspeyres vs Paasche}
\small
\begin{center}
\begin{tabular}{p{0.22\textwidth}p{0.33\textwidth}p{0.33\textwidth}}
\toprule
Criterio & Laspeyres & Paasche \\
\midrule
Ponderación & base & actual \\
Canasta & fija en período base & actual o móvil \\
Lectura & costo de la canasta base hoy & costo de la canasta actual vs base \\
Riesgo & puede ignorar sustitución & puede cambiar la referencia implícita \\
\bottomrule
\end{tabular}
\end{center}
\diferencia{Laspeyres responde con la estructura base; Paasche responde con la estructura actual.}
\end{frame}

\section{Relación entre índices}

\begin{frame}{Relación precio, cantidad y valor}
\[
V=\sum PQ
\]
\begin{itemize}
  \item El valor monetario cambia por dos fuerzas: precios y cantidades.
  \item En el caso simple se cumple exactamente $I(V)=I(P)I(Q)$.
  \item En índices agregados, la relación depende de que las fórmulas de precio y cantidad sean compatibles.
  \item \pregunta{¿El cambio en ventas monetarias se explica por precios, cantidades o ambos?}
\end{itemize}
\end{frame}

\begin{frame}{Identidad fundamental de valor}
\[
I_V=\frac{\sum P_tQ_t}{\sum P_0Q_0}
\]
\begin{block}{Significado económico}
Compara el valor total actual con el valor total del período base.
\end{block}
\begin{itemize}
  \item Si $I_V=1{,}40$, el valor total aumentó $40\%$.
  \item No permite separar por sí solo cuánto se debe a precio y cuánto a cantidad.
  \item \errorcomun{atribuir todo el aumento de valor a inflación sin revisar cantidades.}
\end{itemize}
\end{frame}

\section{Índices de Fisher}

\begin{frame}{Motivación del índice de Fisher}
\begin{itemize}
  \item Laspeyres usa ponderaciones base y Paasche usa ponderaciones actuales.
  \item Cada uno puede entregar resultados distintos cuando cambian los patrones de consumo o producción.
  \item Fisher busca un índice equilibrado combinando ambos.
  \item \pregunta{¿Existe una medida intermedia entre Laspeyres y Paasche?}
\end{itemize}
\end{frame}

\begin{frame}{Índice de precios de Fisher}
\[
P_F=\sqrt{P_L\cdot P_P}
\]
\begin{itemize}
  \item Es la media geométrica entre el índice de precios de Laspeyres y el de Paasche.
  \item \intuicion{equilibra la mirada de canasta base y canasta actual.}
  \item \textbf{Ventaja:} reduce la dependencia de una única estructura de ponderaciones.
  \item \errorcomun{calcular Fisher con media aritmética en vez de media geométrica.}
\end{itemize}
\end{frame}

\begin{frame}{Índice de cantidades de Fisher}
\[
Q_F=\sqrt{Q_L\cdot Q_P}
\]
\begin{itemize}
  \item Combina el índice de cantidades de Laspeyres y el de Paasche.
  \item \pregunta{¿Cuál es el cambio físico equilibrando precios base y actuales?}
  \item Es útil cuando ninguna estructura de precios parece suficiente por sí sola.
\end{itemize}
\end{frame}

\begin{frame}{Comparación final: Laspeyres, Paasche y Fisher}
\small
\begin{center}
\begin{tabular}{p{0.20\textwidth}p{0.27\textwidth}p{0.27\textwidth}p{0.16\textwidth}}
\toprule
Índice & Pondera con & Fortalezas & Cuidado \\
\midrule
Laspeyres & período base & simple, comparable, canasta fija & sustitución \\
Paasche & período actual & refleja estructura actual & base implícita móvil \\
Fisher & media geométrica & equilibrio entre ambos & requiere calcular dos índices \\
\bottomrule
\end{tabular}
\end{center}
\begin{block}{Regla práctica}
Use Laspeyres para seguimiento con canasta fija, Paasche para estructura actual y Fisher como síntesis equilibrada.
\end{block}
\end{frame}

\section{Ejemplo integrador}

\begin{frame}{Ejemplo integrador: datos}
\small
\begin{center}
\begin{tabular}{lrrrr}
\toprule
Producto & $P_0$ & $Q_0$ & $P_t$ & $Q_t$ \\
\midrule
A & 10 & 100 & 12 & 90 \\
B & 20 & 50 & 26 & 60 \\
C & 5 & 200 & 6 & 240 \\
\bottomrule
\end{tabular}
\end{center}
\begin{itemize}
  \item $\sum P_0Q_0=3000$.
  \item $\sum P_tQ_0=3700$.
  \item $\sum P_0Q_t=3300$.
  \item $\sum P_tQ_t=4080$.
\end{itemize}
\end{frame}

\begin{frame}{Ejemplo integrador: cálculo e interpretación}
\small
\[
P_L=\frac{3700}{3000}=1{,}233 \qquad
P_P=\frac{4080}{3300}=1{,}236 \qquad
P_F=\sqrt{1{,}233\cdot1{,}236}=1{,}235
\]
\[
I_V=\frac{4080}{3000}=1{,}36
\]
\begin{itemize}
  \item Los precios aumentaron aproximadamente $23{,}5\%$ según Fisher.
  \item El valor total aumentó $36\%$.
  \item Como el valor crece más que los precios, también hubo efecto de cantidades.
  \item \errorcomun{decir que el valor subió $36\%$ solo por inflación.}
\end{itemize}
\end{frame}

\section{Errores comunes}

\begin{frame}{Errores comunes en números índices}
\small
\begin{enumerate}
  \item Confundir índice con porcentaje.
  \item Interpretar $1{,}20$ como aumento de $1{,}20\%$ en vez de $20\%$.
  \item Cambiar inadvertidamente el período base.
  \item Confundir precio con valor.
  \item Confundir cantidad con valor.
  \item Aplicar Laspeyres usando ponderaciones actuales.
  \item Aplicar Paasche usando ponderaciones base.
  \item Olvidar que Fisher es una media geométrica.
  \item Interpretar causalidad económica a partir de un índice descriptivo.
\end{enumerate}
\end{frame}

\begin{frame}{Checklist de la Unidad 3}
\begin{enumerate}
  \item ¿Identifico claramente período base y período actual?
  \item ¿Distingo índice, variación porcentual e índice base 100?
  \item ¿Separo precio, cantidad y valor?
  \item ¿Reconozco si el índice es simple, complejo, ponderado o no ponderado?
  \item ¿Uso ponderaciones base para Laspeyres y actuales para Paasche?
  \item ¿Calculo Fisher como media geométrica?
  \item ¿Interpreto el resultado sin inventar causalidad?
\end{enumerate}
\end{frame}

\begin{frame}{Cierre de la unidad}
\centering
\Large
Los números índices convierten datos temporales en comparaciones relativas claras.
\vspace{0.8em}

\normalsize
La interpretación correcta depende de tres decisiones: base, ponderación y magnitud analizada.
\end{frame}

\end{document}
